\documentclass{beamer}
\usepackage{beamerthemeshadow}
\usepackage{graphicx}
\usepackage{color}
\usepackage[utf8]{inputenc}
\usepackage{hyperref}
\usepackage[serbian]{babel}
\usepackage[flushleft]{threeparttable}
	\definecolor{chocolate}{rgb}{0.82, 0.41, 0.12}
\setbeamercolor{structure}{fg=chocolate}
\newcommand{\frametitlesection}{
    \insertsectionhead{} - \insertsubsectionhead{}
}

\def\d{{\fontencoding{T1}\selectfont\dj}}
\def\D{{\fontencoding{T1}\selectfont\DJ}}

\title{ChatGPT i njegov uticaj na opšte znanje}
\author{Matija Stanković}
\institute{Matematički fakultet \\ Univerzitet u Beogradu}
\date{Februar 2026}

\begin{document}

\begin{frame}
  \titlepage
\end{frame}

\begin{frame}{Sadržaj}
  \tableofcontents[hideallsubsections]
\end{frame}


\section{Uvod}

\subsection{Veštačka inteligencija - nekada i danas}

\begin{frame}{\frametitlesection}
Veštačka inteligencija nekad
\begin{itemize}
    \item Nastaje sredinom XX veka
    \item \textbf{Ideja: mašine koje imitiraju ljudsko razmišljanje}
    \item Prvi sistemi zasnovani na pravilima i simboličkom zaključivanju
    \item Ograničena primena i mali domet
\end{itemize}

\vspace{0.3cm}

Veštačka inteligencija danas
\begin{itemize}
    \item Razvoj mašinskog učenja i obrade velikih količina podataka
    \item Veštačka inteligencija dostupna širokom krugu korisnika
    \item Prisutna u svakodnevnom životu (pretraživači, društvene mreže, obrazovanje, poslovanje...)
\end{itemize}
\end{frame}

\subsection{Prelomna tačka razvoja - LLM modeli}

\begin{frame}{\frametitlesection}
\textbf{Veliki jezički modeli (LLM)}
\begin{itemize}
    \item Trenirani na ogromnim količinama tekstualnih podataka
    \item Omogućavaju:
    \begin{itemize}
        \item razumevanje prirodnog jezika
        \item generisanje smislenih odgovora
        \item komunikaciju čovek računar
    \end{itemize}
\end{itemize}

\vspace{0.3cm}

Primeri najpoznatijih LLM modela
\begin{itemize}
    \item ChatGPT (OpenAI)
    \item GPT-4 (OpenAI)
    \item Gemini (Google)
    \item Claude (Anthropic)
    \item LLaMA (Meta)
\end{itemize}
\end{frame}

\section{LLM modeli i ChatGPT}

\subsection{Pregled tehnologije}

\begin{frame}{\frametitlesection}

\begin{itemize}
    \item Razvoj LLM modela zasnovan je na \textbf{dubokom učenju}
    \item Omogućavaju značajan napredak u \textbf{obradi prirodnog jezika}
    \item Trenirani na velikim količinama tekstualnih podataka
\end{itemize}

\vspace{0.3cm}

\begin{itemize}
    \item Ključni tehnološki koncept:
    \begin{itemize}
        \item \textbf{Transformer arhitektura} (2017)
        \item Obrada čitavih sekvenci teksta odjednom
        \item Uočavanje odnosa između reči u različitim delovima teksta
    \end{itemize}
\end{itemize}

\vspace{0.3cm}

\begin{itemize}
    \item Za razliku od ranijih sistema:
    \begin{itemize}
        \item nisu zasnovani na ručno definisanim pravilima
        \item poseduju sposobnost \textbf{generalizacije}
    \end{itemize}
\end{itemize}

\end{frame}

\subsection{ChatGPT kao primer LLM modela}

\begin{frame}{\frametitlesection}

\begin{itemize}
    \item Konverzacijski sistem zasnovan na \textbf{GPT arhitekturi}
    \item Razvijen od strane kompanije \textbf{OpenAI}
    \item Komunikacija sa korisnicima putem prirodnog jezika
\end{itemize}

\vspace{0.3cm}

\begin{itemize}
    \item Osnovne karakteristike:
    \begin{itemize}
        \item razumevanje pitanja izraženih svakodnevnim jezikom
        \item generisanje koherentnih i kontekstualnih odgovora
        \item prilagođavanje odgovora toku razgovora
        \item široka primena (obrazovanje, programiranje, informisanje)
    \end{itemize}
\end{itemize}

\vspace{0.3cm}

\begin{itemize}
    \item Ključna prednost: \textbf{visoka dostupnost i jednostavnost korišćenja}
    \end{itemize}

\end{frame}

\subsection{Ograničenja i izazovi LLM modela}

\begin{frame}{\frametitlesection}

\begin{itemize}
    \item \textbf{Nepouzdanost odgovora}
    \begin{itemize}
        \item mogućnost generisanja netačnih ili izmišljenih informacija
    \end{itemize}

    \item \textbf{Odsustvo stvarnog razumevanja}
    \begin{itemize}
        \item modeli nemaju svest niti ljudsko razumevanje značenja
        \item odgovori se zasnivaju na statističkim obrascima
    \end{itemize}

    \item \textbf{Pristrasnost u podacima}
    \begin{itemize}
        \item pristrasnosti iz trening podataka mogu se preneti na odgovore
    \end{itemize}

    \item \textbf{Zavisnost korisnika od modela}
    \begin{itemize}
        \item smanjenje samostalnog razmišljanja i provere informacija
    \end{itemize}
\end{itemize}

\vspace{0.3cm}

\textit{Posebno izraženo u obrazovanju - ChatGPT treba posmatrati kao alat za podršku, a ne zamenu za razumevanje i kritičko razmišljanje.}

\end{frame}


\section{ChatGPT kao izvor informacija}
\subsection{Pristup znanju i dostupnost informacija}

\begin{frame}{\frametitlesection}

\begin{itemize}
    \item \textbf{Laka dostupnost}
    \begin{itemize}
        \item dostupan putem interneta
        \item jednostavan za korišćenje bez tehničkog predznanja
    \end{itemize}

    \item \textbf{Brz pristup informacijama}
    \begin{itemize}
        \item nema potrebe za dugotrajnim pretraživanjem izvora
        \item odgovori prilagođeni pitanju i nivou korisnika
    \end{itemize}

    \item \textbf{Prednosti uz pravilnu upotrebu}
    \begin{itemize}
        \item olakšava učenje i razumevanje pojmova
        \item koristan u samostalnom i neformalnom učenju
    \end{itemize}

    \item \textbf{Rizici nekritičke upotrebe}
    \begin{itemize}
        \item prihvatanje netačnih informacija
        \item površno usvajanje znanja
    \end{itemize}
\end{itemize}

\end{frame}


\subsection{Tačnost odgovora i problem halucinacija}

\begin{frame}{\frametitlesection}

\begin{itemize}
    \item \textbf{Uverljivi, ali ne uvek tačni odgovori}
    \begin{itemize}
        \item gramatički ispravni i logično formulisani
        \item tačnost nije garantovana
    \end{itemize}

    \item \textbf{Halucinacije}
    \begin{itemize}
        \item generisanje informacija koje zvuče uverljivo, ali nisu zasnovane na činjenicama
    \end{itemize}

    \item \textbf{Najčešći problemi}
    \begin{itemize}
        \item izmišljanje činjenica, izvora ili događaja
        \item pogrešno povezivanje tačnih informacija
        \item nemogućnost provere izvora u realnom vremenu
    \end{itemize}

    \item \textbf{Zaključak}
    \begin{itemize}
        \item ChatGPT nije autoritativni izvor znanja
        \item neophodna dodatna provera informacija
    \end{itemize}
\end{itemize}

\end{frame}

\subsection{Uticaj na način pretrage i učenja}

\begin{frame}{\frametitlesection}

\textbf{Promena načina učenja i pretrage}
\begin{itemize}
    \item Direktno postavljanje pitanja umesto klasične internet pretrage
    \item Očekivanje brzih i konkretnih odgovora
\end{itemize}

\vspace{0.3cm}
\textbf{Značaj formulacije upita (prompt-a)}
\begin{itemize}
    \item Kvalitet odgovora zavisi od preciznosti pitanja
    \item Izostavljanje konteksta $\rightarrow$ nepotpuni ili pogrešni odgovori
\end{itemize}

\vspace{0.3cm}
\textbf{Prednosti i rizici}
\begin{itemize}
    \item \textbf{Prednosti:} ušteda vremena, personalizovano i samostalno učenje
    \item \textbf{Rizici:} smanjenje kritičkog razmišljanja, iluzija razumevanja,
    bezbednosni i privatnosni problemi
\end{itemize}

\end{frame}


\section{Uticaj ChatGPT-a na opšte znanje}
\subsection{Promena uloge opšteg znanja}

\begin{frame}{\frametitlesection}
\begin{itemize}
    \item \textbf{Opšte znanje} - skup osnovnih informacija korišćenih u svakodnevnom životu
    \item Tradicionalno sticano kroz:
    \begin{itemize}
        \item formalno obrazovanje
        \item knjige, medije i lično iskustvo
    \end{itemize}
    \item Sa pojavom ChatGPT-a:
    \begin{itemize}
        \item znanje postaje dostupno „na zahtev“
        \item \textbf{fokus se pomera sa šta znamo na kako postavljamo pitanje}
    \end{itemize}
    \item Opšte znanje kao:
    \begin{itemize}
        \item veština snalaženja u informacijama
        \item manje oslanjanje skup trajno zapamćenih činjenica
    \end{itemize}
\end{itemize}
\end{frame}


\subsection{Pojednostavljivanje i percepcija znanja}

\begin{frame}{\frametitlesection}
\begin{itemize}
    \item ChatGPT pojednostavljuje složene teme
    \item Omogućava:
    \begin{itemize}
        \item brz i pristupačan \textbf{osnovni uvid}
        \item upoznavanje sa novim oblastima
    \end{itemize}
    \item Pozitivan efekat:
    \begin{itemize}
        \item veća informisanost korisnika
        \item niži prag za ulazak u novu temu
    \end{itemize}
    \item Potencijalni rizik:
    \begin{itemize}
        \item \textbf{iluzija razumevanja}
        \item izostavljanje složenijih aspekata problema
        \item površno znanje bez dubljeg uvida
    \end{itemize}
\end{itemize}
\end{frame}


\subsection{Opšte znanje i akademski kontekst}

\begin{frame}{\frametitlesection}
\begin{itemize}
    \item ChatGPT se često koristi za:
    \begin{itemize}
        \item opšta pitanja
        \item kratka objašnjenja pojmova
        \item svakodnevne nedoumice
    \end{itemize}
    \item Sve češća primena u:
    \begin{itemize}
        \item obrazovanju
        \item akademskoj zajednici
        \item pisanju i istraživanju
    \end{itemize}
    \item U akademskom kontekstu:
    \begin{itemize}
        \item veći uticaj na proces učenja
        \item veća odgovornost korisnika
        \item ozbiljnije posledice nepravilne upotrebe
    \end{itemize}
\end{itemize}
\end{frame}


\section{ChatGPT u obrazovanju}

\subsection{Upotreba ChatGPT-a u nastavi}

\begin{frame}{\frametitlesection}
\begin{itemize}
    \item Sve češća primena u visokom obrazovanju
    \item Najčešće se koristi za:
    \begin{itemize}
        \item objašnjavanje složenih pojmova
        \item razumevanje gradiva
        \item generisanje primera
        \item pomoć pri pisanju i strukturiranju tekstova
    \end{itemize}
    \item Omogućava \textbf{personalizaciju učenja}:
    \begin{itemize}
        \item prilagođavanje nivou znanja studenta
        \item fleksibilan tempo učenja
    \end{itemize}
    \item Važno ograničenje:
    \begin{itemize}
        \item nije zamena za nastavnika i samostalno razmišljanje
    \end{itemize}
\end{itemize}
\end{frame}


\subsection{Iskustva i rezultati istraživanja}

\begin{frame}{\frametitlesection}
\begin{itemize}
    \item Studenti brzo prihvataju ChatGPT i integrišu ga u svakodnevno učenje
    \item Najčešće se koristi za:
    \begin{itemize}
        \item objašnjavanje gradiva
        \item proveru razumevanja
        \item dobijanje ideja za dalji rad
    \end{itemize}
    \item Ređe se koristi kao konačni izvor informacija
    \item Uočeni problemi:
    \begin{itemize}
        \item precenjivanje tačnosti odgovora
        \item izostanak provere informacija
        \item rešavanje zadataka bez razumevanja postupka
    \end{itemize}
    \item Zaključak istraživanja:
    \begin{itemize}
        \item potreba za edukacijom o odgovornoj upotrebi
        \item fokus na kritički odnos, a ne zabranu
    \end{itemize}
\end{itemize}
\end{frame}


\subsection{Etika i akademski integritet}

\begin{frame}{\frametitlesection}
\begin{itemize}
    \item \textbf{Glavni problem:}
    \begin{itemize}
        \item zloupotreba pri izradi domaćih zadataka i seminarskih radova
        \item nejasna granica između pomoći i neetičkog ponašanja
    \end{itemize}
    \item \textbf{Akademski integritet:}
    \begin{itemize}
        \item samostalnost u radu
        \item poštovanje autorstva
        \item transparentno korišćenje izvora
    \end{itemize}
    \item Nekritičko korišćenje ChatGPT-a može dovesti do:
    \begin{itemize}
        \item prikrivenog plagijata
        \item narušavanja akademskog integriteta
    \end{itemize}
    \item Rešenje:
    \begin{itemize}
        \item jasna pravila upotrebe
        \item ChatGPT kao pomoćni alat, a ne zamena za učenje
    \end{itemize}
\end{itemize}
\end{frame}


\subsection{Mogućnosti i ograničenja primene ChatGPT-a u obrazovanju}

\begin{frame}{\frametitlesection}
\begin{itemize}
    \item ChatGPT može da unapredi obrazovni proces
    \item Nije zamena za učenje, već:
    \begin{itemize}
        \item alat za podršku učenju i pomoć pri razumevanju i vežbanju gradiva
    \end{itemize}
    \item \textbf{Primer pravilne upotrebe:}
    \begin{itemize}
        \item generisanje primera zadataka za samostalnu analizu
    \end{itemize}
    \item \textbf{Uloga nastavnika:}
    \begin{itemize}
        \item korišćenje ChatGPT-a kao pomoćnog alata uz obaveznu stručnu kontrolu i nadzor
    \end{itemize}
    \item \textbf{Dugoročni pravac razvoja:}
    \begin{itemize}
        \item namenski AI alati za edukaciju
        \item modeli trenirani na relevantnim nastavnim materijalima
    \end{itemize}
    \item Ključ uspeha: \textbf{edukacija o odgovornoj i pravilnoj upotrebi}
\end{itemize}
\end{frame}


\section{Zaključak}

\begin{frame}{Zaključak}
\begin{itemize}
    \item ChatGPT i drugi LLM modeli predstavljaju sadašnjost i budućnost sticanja znanja
    \item \textbf{Nepravilna upotreba}:
    \begin{itemize}
        \item nekritičko prihvatanje odgovora
        \item oslanjanje na ChatGPT kao jedini izvor znanja
        \item površno razumevanje i širenje netačnih informacija
    \end{itemize}
    \item \textbf{Pravilna upotreba}:
    \begin{itemize}
        \item pomoć pri razumevanju i učenju
        \item usmeravanje daljeg istraživanja
        \item ubrzavanje procesa sticanja znanja
    \end{itemize}
\end{itemize}
\textbf{Da li ChatGPT utiče pozitivno ili negativno na opšte znanje?}
\end{frame}


\begin{frame}{Literatura}
\footnotesize
\begin{thebibliography}{99}

\bibitem{russell2020}
S. Russell, P. Norvig, \textit{Artificial Intelligence: A Modern Approach}, Pearson, 2020.

\bibitem{bommasani2021}
R. Bommasani et al., \textit{On the Opportunities and Risks of Foundation Models}, arXiv:2108.07258, 2021.

\bibitem{bender2021}
E. M. Bender et al., \textit{On the Dangers of Stochastic Parrots}, ACM FAccT, 2021.

\bibitem{openai2023}
OpenAI, \textit{GPT-4 Technical Report}, \url{https://arxiv.org/abs/2303.08774}

\bibitem{seminarMilos}
M. Pavlović, \textit{Primena veštačke inteligencije u obrazovanju}, Simpozijum \textit{Matematika i primene}, 2025.

\end{thebibliography}
\end{frame}

\begin{frame}{}
\centering
\vspace{1.5cm}

{\Large \textbf{Hvala na pažnji!}}

\vspace{1cm}

{\large Pitanja?}

\end{frame}

\end{document}


\end{document}